% LAB BOREDOM — module L2, The Apparatus and the Data. GUIDED DRAFTING (desktop recipe).
% Session 2026-08-06, continuing the Section 2 run. Spec: outline v7 §L2. Rulings: register v8.
% Protocol: Stage-3 guided batches, approval gate each batch; bare ¶-numbers when presenting.
% Sources for THIS module differ from L1: there is no L1-style pack. Evidence comes from
%   (a) boredom-analysis-v5-2026-08-05/ (METHODS.md, FINDINGS.md, out/*.csv) — the ONLY tree;
%   (b) PROSE-METHODS-CRITERION-VALIDITY-v2 and PROSE-METHODS-GAZE-BASELINE-v2 — DO NOT REWRITE;
%   (c) our own published papers, quoted only through gated spans.
% Standing rulings carried from L1 (see L1-DRAFT-v1.tex header for the full list):
%   first+last name at first mention · NO vague "instrument" · American spelling · MLA ·
%   no bold run-in heads · "demonstrated" never "prove" · all three films treated equally.
%
% ★ COHORT STRUCTURE — CORRECTED AND AUTHOR-CONFIRMED 2026-08-06. This overturns a standing
%   finding and the outline's L2.1 framing. Record the change in register v9.
%   (1) The 2023 paper (P1) reported THREE subjects. Verified in the document four ways:
%       the Methods sentence "three individuals, two of whom had no prior experience...";
%       Table 1's three demographic rows (ages 20, 21, 33); results discussing Subject 1-3
%       only; the limitations naming Subject 2's artifacts and Subject 3's impedances.
%       Four copies of the PDF across four repos are byte-identical (md5 dc8c888f...).
%       AUTHOR CONFIRMED 2026-08-06 that three is correct for the published paper.
%   (2) Round 1 as a COLLECTION ROUND has eight subjects (S01-S08). The 2023 paper published
%       three of them. These are different facts about different things.
%   (3) ★ F-11 IS WITHDRAWN (author ruling 2026-08-06). The register asserts "the published
%       gaze cohort cannot contain Round 2," inferred from the 2024 paper's reference access
%       dates (all Jan 16, 2024) against Round 2 being recorded that March. The author, who
%       ran the study, confirms the 2024 paper's twelve IS Round 1's eight plus Round 2's
%       four, using the SAME Round 1 data. Arithmetic corroborates: only twelve people exist
%       in the dataset. Citation access dates date the references, not the submission.
%   (4) CONSEQUENCE FOR L2.1: outline v7's "the two twelves are different twelves" trap is
%       GONE. Replaced by continuity — the published gaze cohort IS this analysis's cohort.
%   (5) Participants include COMPUTER SCIENCE undergraduates (author, 2026-08-06) alongside
%       the BME undergraduates, MEng BME students, and first-/second-year medical students
%       the 2024 paper's own participant paragraph names.
%   (6) IRB: use the 2024 form — UCI IRB Exempt No. 20232678. (The 2023 paper renders it
%       "UCI IRB exempt No. 2678".)
%
% ★ PARHI AND AYINALA RE-PINNED: 2014, NOT 2013. IEEE Transactions on Circuits and Systems I,
%   vol. 61, no. 1, January 2014. Outline v7 §L2.4 flagged the suspicion; confirmed from the
%   PDF front matter. Works cited takes 2014.
%
% O-14: author ruling 2026-08-06 — multiplicity in the criterion family is handled as a
%   FOOTNOTE, not a computed correction. Revisable later.
% O-15: author ruling 2026-08-06 — DEFERRED to the point in the draft where it lands (L2.8).
%   The ****** stays until then.
% ============================================================================
% L2 REVISION BATCH (2026-08-21, author-ruled during appendix drafting; backup
% .backups/L2-revision-2026-08-21/):
%   L2-E1 administration corrected per author statement: SIX questions in the break; the
%     fatigue question ("how tired right now") asked JUST BEFORE each film began.
%   L2-E2 "after each film" -> "for each film"; appendix footnote notes the ADMINISTERED
%     wording (D-H: administered text is canonical; ****** designation stays provisional
%     per the renumbering rule).
%   L2-E3 parsing passage updated to the D-39 counts (239/5/1/7/0, none missing; six
%     never-bored minutes cells -> the n=30 explanation; one NA->scale-floor; the flagged
%     annotation case). The pre-D-39 "five semantic / two absent" text is superseded.
%   L2-E4 the O-15 ****** marker DELETED — O-15 RESOLVED by author statement 2026-08-21
%     (administration description now lives in E1 and Appendix B).
% ============================================================================
%
% BATCH 1 = L2.1 + L2.2 (drafted 2026-08-06, awaiting review).
%   Survey figures verified in boredom-analysis-v5 METHODS.md §6: seven items =
%   fatigue-prior, boredom (1-9), engagement (1-9), minutes-until-bored, sleep-fight,
%   felt duration, fatigue-after. Derived: felt_duration_ratio = felt/17;
%   depletion = fatigue_after - fatigue_prior. Parsing confidence: exact 239, converted 5,
%   flagged 1, semantic 5, missing 2; every interpreted value written to EXCLUSIONS.log with
%   its raw string; 'Halfway' read as 8.5 min (author-approved 2026-07-30).

% BATCH 1 REV 1 (2026-08-06, author notes):
%   P1: the publication-reconciliation paragraph COLLAPSED INTO A FOOTNOTE at first mention of
%     the papers, since an appendix will carry all publications. The body now narrates the
%     experiment in the order it was run.
%   P2.1: pedantic opener replaced.
%   P2.2: EXACT PROCEDURE ORDER added — calibration, then three 17-minute films with ~5-minute
%     breaks, survey administered in each break. Sources: P1 ("Each video was 17 minutes long
%     and participants were given roughly a 5-minute break in between each video to allow them
%     to restabilize"; headset worn "over their EEG cap") and P2 ("asked to first perform an
%     eye-tracking calibration session, followed by watching engaging, clinical immersion, and
%     boring VR videos").
%   P2.3: SURVEY QUESTIONS pulled verbatim from the 2023 paper's appendix tables and pointed to
%     a new appendix via footnote. SIX of the seven verify verbatim:
%       (1) "On a scale from 1-9 how bored were you while watching this video? (1=Not bored at
%           all, 9= Extreme boredom)"
%       (2) "On a scale from 1-9 how interested were you while watching this video? (1=Not
%           engaged at all, 9=Intense engagement)"  [Table 4 variant: "how interested and
%           engaged in the video were you?"]
%       (3) "Roughly how many minutes into the video until you started to get bored?"
%       (4) "Did you fall asleep/fight to stay awake? (1=Completely awake, 9=Fell asleep)"
%       (5) "Relatively speaking, how long did the video feel to you?"
%       (6) "On a scale from 1-9 how tired were you before conducting the experiment
%           (1=completely awake, 9=completely exhausted)"
%     ★ THE SEVENTH ITEM (fatigue AFTER the episode, `fatigue_after` in the v5 tree, which
%       `depletion` is computed from) DOES NOT APPEAR in the 2023 paper's printed tables.
%       Its wording is ****** UNVERIFIED and must be recovered from the raw survey files
%       before the appendix is assembled. Do not invent it.
%   P2.4/2.5: the parsing passage REWRITTEN. Key clarification the author asked for: the counts
%     sum to 252 = 36 episodes x 7 items, i.e. every answer collected. Verified: 239+5+1+5+2=252.

Twelve people took part, eight in the first round of recording and four in the second, all
recruited at the University of California, Irvine from undergraduate biomedical engineering,
master's biomedical engineering, computer science, and first- and second-year medical
students, under IRB exemption 20232678.\footnote{Both published studies draw on these same
twelve. The 2023 paper reported three of the first round's eight participants; the 2024 paper
reported all twelve, using for the first-round participants the same recordings the earlier
paper had drawn on. The analysis presented here pools all twelve. Both papers are reproduced
in Appendix ******.} Each participant was fitted with an EEG cap and an HP Reverb G2 Omnicept
headset worn over it, then completed an eye-tracking calibration in which they fixed on points
around a circle at the center of the display and registered each with a trigger pull on the
headset's controller. They then
watched three films of seventeen minutes each---a boring film, an interesting film, and
clinical immersion footage---while EEG, eye tracking, and the headset's other sensors recorded
continuously. Between films came a break of roughly five minutes, long enough for the
participant to settle before the next recording began. Six of the survey's seven questions
were administered in that break, their answers referring to the film just watched rather than
to the session as a whole; the remaining question---how tired the participant felt right
then---was asked just before each film began. Thirty-six episodes result: twelve participants
by three films.

The survey asked seven questions for each film.\footnote{The full survey, in its administered
wording, is reproduced in Appendix ******.} Two of them do the work the rest of this section depends on. Boredom and
engagement were asked separately, each on its own nine-point scale, rather than as two ends of a
single scale running from bored to engaged. That is why disagreement inside a participant's own
answers is visible at all: someone can rate boredom high and engagement high in the same breath,
and some did. The remaining questions asked how many minutes passed before boredom set in, how
hard the participant fought to stay awake, how long the film felt, and how tired they were
before and after watching. Two further quantities are calculated from these rather than asked
directly. Felt duration becomes a ratio against the film's actual seventeen minutes, so a film
that felt like half an hour scores near 1.8 and one that passed quickly scores below 1.
Depletion is simply the second fatigue rating minus the first: how much more tired a participant
was when a film ended than when it began, which separates the film's own effect on alertness
from whatever tiredness they brought into the session.

The questions were put to participants aloud and their answers written down by the proctor, so
the surveys survive as free text rather than as selections from fixed options, and turning that
text into numbers required reading each answer. Those readings are recorded rather than made
quietly. Seven questions across thirty-six episodes gives 252 answers, and none is missing.
Of these, 239 were plain numbers within the expected range and needed no interpretation at
all. Five required a unit or a range to be converted---``25 sec'' became 0.42 minutes,
``10--12'' became 11. One arrived in range but carrying an annotation, which is logged with
it. Seven were not numbers yet carried their meaning plainly: six answers to how many minutes
passed before boredom set in record that the participant never became bored---a finding
rather than a missing value, kept as its own category rather than filled in or counted as
zero, which is why that item is tested on thirty values rather than thirty-six---and one
``NA'' on a 1--9 item means the film produced none of the asked-after quality and is read as
1, the scale's floor. Every
interpreted value is stored alongside the exact words the participant wrote, so any reading can
be checked against its original.\footnote{One answer reading ``Halfway'' was read as 8.5
minutes, half of the seventeen-minute film.} What the survey can register follows from the
definition it was built under. Fahlman's account supplies that definition, not the wording of
the questions, which put it to participants in ordinary language---and that is what makes these
seven items a first-form measurement, able to record the boredom a person notices and can
attribute, and unable in principle to record the rest.

% ============================================================================
% BATCH 2 = the criterion method (outline v7 §L2.8, now the section's ONLY method).
% Drafted 2026-08-07 under the author's ruling of the same date:
%   ★ FILM-SEPARATION STATISTICS ARE EXCLUDED FROM THE DISSERTATION ENTIRELY.
%     Not demoted, not appendixed. The 99 pairwise contrasts, the Friedman/Wilcoxon
%     across-film tests, and the nine Holm-surviving contrasts do not appear.
%     Boredom is a response, not a property of a film, so the organizing variable is
%     what each participant reported about each episode.
%   ★ The criterion analysis IS the experiment.
%   ★ Body carries the Spearman/Pearson comparison chart; the fuller criterion table
%     for non-load-bearing channels goes to the appendix.
%   ★ Sleep-fight: within-participant result in the body as a convergent finding, with
%     the correlation beside it carrying the shared-method caveat.
%   ★ METHOD MUST BE EXPLAINED FOR A READER WITHOUT STATISTICAL TRAINING — not only the
%     results but what was computed and why that computation and not another.
% CONSEQUENCE: PROSE-METHODS-CRITERION-VALIDITY-v2's "do not rewrite" status LAPSES for
%   register. Its numbers and logic are preserved; its pitch is not, because v2 assumes a
%   statistically literate reader and the author has ruled otherwise. Numbers below are
%   read from boredom-analysis-v5 out/criterion-validity.csv after the D-39 re-run.
% NEW ANALYSIS (2026-08-07, not in outline v7): sleep-fight against reported boredom and
%   engagement, and the within-participant most-boring vs most-engaging paired test.
%   Computed this session; verified figures in the batch-2 ledger at the foot of this file.
% ============================================================================

Everything that follows compares two things: what a participant's body did during an episode
and what that participant said about the same episode afterward. Four decisions govern how that
comparison was made, and each exists to remove a way the comparison could mislead. They are, in
the order taken up below, within-subject standardization, the parallel reporting of rank-order
and product-moment correlation, a within-cluster Monte Carlo permutation test in place of the analytic one,
and a paired within-subject contrast evaluated by an exact sign test.

The first decision is to measure every participant against themselves, which is done by
converting each value to a within-subject standard score, or z-score. People differ enormously
in ways that have nothing to do with boredom. One person's eyes range widely across any scene;
another's barely move. One person uses the whole nine-point scale; another never says more than
seven about anything. Pooling the thirty-six episodes and correlating them directly would let
those standing differences between people dominate the result, so that much of what came out
would describe who has a restless eye or a generous hand with a rating scale rather than
anything about boredom. Standardizing within each participant removes it: every value becomes
its distance from that participant's own mean, divided by that participant's own standard
deviation. In plain terms, each person is scored against their own average and their own
typical spread, so a participant whose gaze wandered furthest on the film they found dullest
registers identically whether their eyes are naturally still or naturally restless. What each
person contributes is the shape of their own three episodes and nothing else, which is precisely
the quantity the question is about.

The second decision is to report the association two ways, as Spearman's rank-order correlation
and as Pearson's product-moment correlation. Spearman's coefficient asks only about order: when
a participant's episodes are ranked from least to most boring, do the physiological measurements
fall in the same order? It is computed on ranks rather than on the values themselves, which is
what makes it the appropriate primary statistic here, because it assumes only that the rating
scale is ordinal---that a 9 means more than an 8---and not that it is interval-scaled, that the
distance from a 3 to a 4 is the same size as the distance from an 8 to a 9. For a
self-report scale that stronger assumption cannot be guaranteed. Pearson's coefficient asks
whether the relationship is linear: does each additional point of reported boredom correspond to
roughly the same additional amount of gaze wandering all the way up the scale? It carries the
interval assumption Spearman's avoids, so the conclusions do not rest on it. Both are reported
because their agreement is itself a result. Where the two coefficients converge, the association
is not merely monotonic but close to proportional, which is the condition that would have to
hold before such a measurement could stand in for a rating rather than merely accompany it.

The third decision concerns how an association is judged unlikely to have arisen by chance, and
it replaces the analytic probability with a within-cluster Monte Carlo permutation test. The textbook
formula assumes the thirty-six episodes are thirty-six independent observations. They are not:
they are twelve participants contributing three episodes each, and episodes from one person
resemble one another for reasons unrelated to the films. That clustering violates the
independence assumption and inflates significance, sometimes badly. The Monte Carlo test builds
its null distribution from the data instead. Each participant's three ratings are randomly
reassigned among that participant's own three episodes, the coefficient is recomputed, and the
procedure is repeated twenty thousand times. The reported probability is the proportion of those
twenty thousand rearrangements producing an association at least as strong as the observed one.
Put without the vocabulary: the connection between body and report is deliberately broken twenty
thousand times over, and the question is how often blind chance manages to do as well as the
real data did. Because the reshuffling happens strictly inside each participant, the clustering
is respected rather than assumed away.

The fourth decision applies to the sleep item and answers an objection to it, by way of a paired
within-subject contrast. A participant who arrived short of sleep will fight drowsiness in every
episode, and that has nothing to do with any film; drowsiness therefore carries a component
belonging to the person rather than to what they watched, which in a between-episode comparison
would appear as noise or, worse, as a spurious effect. Pairing removes it. For each participant,
the episode that participant rated most boring is set against the episode the same participant
rated most engaging, and the drowsiness reported in each is compared. Any tiredness carried into
the session appears on both sides of that difference and cancels, so what remains is the effect
of that participant's own experience of the two episodes. Because the design yields one
difference per participant, significance is assessed by an exact sign test---counting how many
participants went in each direction and computing the exact binomial probability of a split at
least that lopsided---supported by an exact sign-flip permutation over the same differences.
With twelve participants the test has a hard floor: even a perfect split cannot produce a
probability below roughly one in a thousand, so the reported value should be read as the
strongest result this sample size can express rather than as a precise measure of rarity.

% ============================================================================
% BATCH 3 = the stimuli, and the two-apparatus problem (outline §L2.3 + §L2.4).
% Drafted 2026-08-07. Verified this session against archon-cli-v3 and the v5 tree:
%   Q-P1a "State boredom often derives from information or environmental stimuli that is
%     monotonous, redundant, and/or meaningless" -- EXACT, King & Salvo P1 (2023).
%   Q-P1b "Due to the general familiarity our demographic has with Microsoft Word along with
%     the age and obsolescence of the interface, we believe this video meets the criteria
%     necessary to induce state boredom" -- EXACT, same source.
%   Rates: R1 median 3.00 Hz (2.79-3.00 after harmonization); R2 median 120.1 Hz. Factor ~40.
%   Poolable: pupil, cognitive load, HR, eye closure/openness, self-report. NOT poolable:
%     gaze_dev_median, gaze_dev_variance (5-point rolling median window = ~42 ms at 120 Hz
%     vs ~1.5 s at 3 Hz); decimating R2 to R1's rate moves variance by 0.65-0.84x (mean 0.74).
%   Blink: R1's median sampling interval ~334 ms = the mean duration of a spontaneous blink,
%     so a blink occupies 0-1 samples and is unresolvable at any threshold. 20 Hz is the floor.
%     Blink rate is therefore a Round 2 measure only.
%   Openness scale harmonization: R1 percentage (0-100) vs R2 fraction (0-1), R1 x 0.01.
% NOTE: film-separation consequences of the rate difference are NOT stated, per the ruling.
%   The rate difference is presented as a constraint on WHICH MEASURES CAN BE POOLED, which is
%   what it is, rather than as a caveat about comparing films.
% ============================================================================

Three films were shown, each seventeen minutes long, and each was chosen before any data
existed to serve a particular role. The boring film is a Microsoft Word tutorial recorded in
1989, selected because boredom has long been associated with material that is redundant,
monotonous, of low intensity, or meaningless---a characterization Fahlman and colleagues draw
together from a literature reaching back to Berlyne, and the one our own study invoked when
choosing the stimulus.\footnote{Fahlman and colleagues present it as a conceptualization the
field has offered rather than as their own definition, citing Berlyne (1960), Fiske and Maddi
(1961), Geiwitz (1966), Hebb (1966), Mikulas and Vodanovich (1993), and Posner, Russell and
Peterson (2005). Our 2023 paper paraphrases their summary.} A software
tutorial for a program the participants already knew, delivered through an interface three
decades obsolete, was expected to supply those qualities at once. The paper states the
reasoning plainly: ``Due to the general familiarity our demographic has with Microsoft Word
along with the age and obsolescence of the interface, we believe this video meets the criteria
necessary to induce state boredom.'' The interesting film covers alleged alien reproduction
vehicles and the deaths of scientists associated with alternative energy research, chosen on
the expectation that a topic of live public fascination would hold attention regardless of
whether its claims withstand scrutiny---accuracy was not the criterion, since the film was
selected to engage rather than to inform. The clinical film is a clip of a spinal deformation
procedure drawn from the platform's own library, and it is the material the platform exists to
deliver. Whether any of the three actually produced the state it was chosen to produce is a
question for the participants rather than for the selection rationale, and each participant
answered it for themselves.

The two rounds were recorded on the same equipment, and what differs between them is the
software that captured what the equipment produced. Participants in both rounds wore the same
headset with the same sensors; the only hardware change is that the EEG cap was dropped for
Round 2. The capture software, however, was still being developed during Round 1 and logged at
a median of three samples per second. By Round 2 it had been brought to the point where every
sensor polled and logged at the full rate of one hundred and twenty. That difference is a
factor of roughly forty, and it governs what can be combined, because measurements differ in
how much they care about how often they were sampled. A
quantity that is an average of instantaneous values---pupil diameter, the proportion of an
episode spent with the eyes closed, heart rate---is unaffected: averaging a hundred and twenty
readings per second and averaging three per second give the same answer, because the underlying
quantity is being sampled rather than constructed, and the sensor was reporting the same thing
in both rounds regardless of how often the software wrote it down. A quantity computed over a
moving window is a different matter. The gaze-deviation measure smooths across five consecutive samples before
anything else happens, and five samples spans about forty-two milliseconds at Round 2's rate
against roughly a second and a half at Round 1's. Those are not the same measurement wearing one
name; recomputing Round 2's gaze variance after thinning it down to Round 1's rate changes the
value by a quarter or more. Every quantity in this analysis therefore carries a flag recording
whether it may be pooled across the two rounds, with the reason attached, and windowed
quantities are computed within a round and never combined across them.

One measurement is lost to the same arithmetic, and stating why is more useful than reporting
it with a caveat. A spontaneous blink lasts about a third of a second. Round 1's median interval
between samples is also about a third of a second, so a blink there occupies either one sample
or none, and no threshold applied to that data can distinguish a blink from a gap in the logged samples.
Around twenty samples per second is the floor for measuring blinks honestly. Blink rate is
consequently a Round 2 quantity only, and no Round 1 blink figure appears anywhere in this
analysis. A smaller mismatch was repairable rather than fatal: the two rounds recorded eye
openness on different scales, one as a percentage and one as a fraction, and the two are brought
onto a single scale before any test that ranks the sizes of differences, since leaving them
mixed would silently give one round's participants forty times the weight of the other's.

% ============================================================================
% BATCH 4 = eye closure and its Round 1 proxy (§L2.5) + the two windows (§L2.6).
% Drafted 2026-08-07. Figures verified in boredom-analysis-v5 METHODS.md this session:
%   R1 has NO continuous openness channel; the per-sample validity flag (validL & validR)
%     is the proxy, aggregated as pct_valid_eye. LICENSED inside Round 2 where both exist:
%     validity proxy vs vendor continuous openness correlate at PEARSON r = 0.998 across
%     the 12 Round 2 cells.
%   R1 closure EPISODES read from the per-eye -1 sentinel in dilL/dilR (parse_window()
%     masks those to NaN -- correct for pupil, fatal for closure). An episode is a maximal
%     run where BOTH eyes read -1.
%   R2 openness VERIFIED BINARY: {0,1} only across 12 cells and 3,416,362 per-eye samples,
%     no intermediate value, so OPENNESS_CLOSED_BELOW = 0.5 is NOT a tuned parameter and the
%     sensitivity sweep over it was redundant and not run. CONSEQUENCE: blink vs extended
%     closure rests ENTIRELY on duration.
%   Blink ceiling 500 ms = VanderWerf mean 334 +/- 67 ms, i.e. mean + ~2.5 SD, rounded;
%     those blinks were recorded while subjects watched video, matching this paradigm.
%     Source pinned: J Neurophysiol 89(5):2784-96, doi:10.1152/jn.00557.2002.
%   Episode counts: R1 4365 (3501 blink-length, 864 extended); R2 5131 (4895 blink, 236 ext).
%   ★ NEW LIMIT (VanderWerf entry, 2026-08-06): the ceiling is a DURATION threshold and blink
%     duration varies systematically with vertical gaze position -- controllable at a fixed
%     screen, NOT in 360 video where gaze is the participant's own behavior. The partition,
%     the episode counts and R2's blink rate all inherit that dependency. Carry to L6.
%   WINDOWS: R1 crop mean 818.0 s (range 661.1-973.7); R2 full mean 1038.9 s (1013.2-1167.8);
%     r1matched = crop minus 120 s each end, mean 798.9 s (773.2-927.8). R1 removed ~2 min
%     inside each 30 s instructed closure as a sync guard (protocol not yet settled), so its
%     window is the middle ~13.6 min of a 17-min film. `full` is primary; both carried
%     through the same tests in window-comparison.csv. THE RESULT BELONGS TO L3, NOT HERE.
% ============================================================================

Eye closure is named for exactly what it measures: the proportion of an episode during which
the participant's eyes were shut. Both rounds record it the same way, and the recording is
unusually clean. The headset reports each eye as open or closed and nothing in between---in
Round 1 the pupil channel writes a sentinel value whenever an eye cannot be seen, and in Round 2
a dedicated per-eye channel carries the same information; across all twelve Round 2 recordings
and more than three million individual eye samples the only values present are zero and one,
with no intermediate value anywhere. The consequence is that the distinction between a blink and
a sustained closure rests entirely on how long the closure lasted, and on nothing else.

An episode counts as closed only where both eyes are shut at once.\footnote{Single-eye readings
occur and are not counted. Averaged across episodes they would add about half a percentage point
of closure time, reaching 2.2 points in the most affected episode. They are excluded because the
data cannot distinguish a genuinely closed eye from an eye the tracker has simply lost---a
headset shifting on the face, an eyelash, a lens, an uneven calibration---and treating a
tracking failure as closure would let how well the equipment happened to fit a particular
participant register as boredom. The conservative rule costs nothing measurable: the two
definitions correlate at 0.998 across the thirty-six episodes, and adopting it moves the
association between closure and reported boredom from +0.460 to +0.475, and between closure and
reported engagement from $-$0.417 to $-$0.435.} Where the two rounds are read from different
columns, they are read to the same rule.

Separating blinks from sustained closures requires a boundary in time, and the boundary is
drawn from VanderWerf and colleagues, who measured spontaneous blinks at an average of 334
milliseconds with a spread of 67, in participants who were themselves watching video, which is
the same situation as this experiment. A closure of half a second or less is therefore treated
as a blink, that being the mean plus about two and a half standard deviations rounded upward,
and anything longer is treated as a sustained closure. The same work established that blink
duration is not a constant: it varies systematically with where the eye is pointed
vertically.\footnote{The eyelid follows the eye. Looking downward brings the upper lid partway
down with it, so a blink begun from a downward gaze starts with the lid already lower and
travels a shorter distance than one begun from a level gaze---the same reflex, differently
sized and differently timed, depending only on where the eye happened to be pointing. The
participants here varied enough for this to matter. Measured against the headset's forward axis,
one Round 2 participant spent five percent of every episode looking more than twenty degrees
below center, and held a median gaze between four and six degrees below center throughout;
another sat closer to level, with a median gaze two to four degrees above center. A blink
measured on the first participant and a blink measured on the second are not being measured
under the same conditions.} That dependency has a particular consequence in this apparatus,
because vertical gaze here is not a controlled variable but a behavior.

Participants were seated and did not turn around, and the screen was fixed in front of them in
every case: the boring and interesting films were conventional flat video, presented as such,
and the clinical footage was a hundred and eighty degrees, in stereoscopic 3D, rather than a
full surround. What was
never fixed is where within that field a participant chose to look. A participant who lets their
gaze fall toward the floor stops seeing the boring or interesting film altogether, and sees only
the floor of the operating room in the clinical footage. That behavior is not incidental to what
this section studies---looking away from a film is one of the more legible things a disengaged
viewer does---and it is also the behavior that changes how a blink measures. The partition into
blinks and sustained closures, the episode counts, and Round 2's blink rate therefore inherit a
dependency on the very conduct under observation, which is a reason to treat the blink
measurements as descriptive and to rest nothing on them.

The two rounds also differ in how much of each film was analyzed, and the difference is not
neutral for a measurement that accumulates. Round 1's recordings were cropped with roughly two
minutes removed inside each of the instructed closures that bracket a film, a sync guard
adopted while the protocol was still being settled, which leaves the middle thirteen and a half
minutes or so of a seventeen-minute film. Round 2's recordings are essentially the whole of it.
Since a participant can only close their eyes for a larger fraction of a window if the window
contains the moments when they did so, a shorter window and a longer one are not exchangeable.
Rather than assume the difference away, both are computed: every quantity is calculated on the
full recording and again on a window trimmed to match Round 1's, and both are carried through
the same tests. The full window is primary. What the comparison between them shows belongs with
the results rather than here.

% ============================================================================
% BATCH 5 = viewing order (§L2.7), the gaze baselines (§L2.9), and what each measurement
% can and cannot say (§L2.10, REBUILT -- its v7 form was organized around separation).
% Drafted 2026-08-07.
%
% ★ CONSEQUENCE OF THE GOVERNING RULING FOR §L2.9: the three-baseline architecture was built
%   to answer TWO questions -- separation and agreement -- with device-forward assigned to
%   separation (where it carried "the only gaze result in this study to survive correction")
%   and per-participant assigned to agreement. With separation excluded, device-forward has
%   no assigned question left. It is therefore reported as a computed alternative rather than
%   as the holder of a question. PROSE-METHODS-GAZE-BASELINE-v2 §0.7-0.8 is superseded on that
%   point; its §0.9 justification for the per-participant baseline is UNAFFECTED and is what
%   the paragraph below renders, because that argument was always from design, not coefficient.
%
% Verified figures (PROSE-METHODS-GAZE-BASELINE-v2 + criterion-validity.csv after D-40):
%   per-participant vs boredom +0.590 (perm 0.0024); vs engagement -0.585 (0.0022)
%   per-file        vs boredom +0.524 (0.0074); vs engagement -0.527 (0.0075)
%   device-forward  vs boredom +0.359 (0.0751); vs engagement -0.391 (0.0515) -- neither survives
%   F-01 axis conventions: R1 Unreal X forward / Y horizontal / Z vertical;
%                          R2 Unity  Z forward / X horizontal / Y vertical (from original MATLAB)
%   Viewing order (F-09/F-14): S01 and S05 ran CLC -> BOR -> INT; the other six Round 1 and all
%     four Round 2 ran INT -> CLC -> BOR. Boring LAST in ten of twelve, second in two. Order
%     fixed after the first two participants; the published paper states the reason.
%     ★ THE TEST OF THE ORDER HYPOTHESIS GOES TO L6, NOT HERE.
%   Appendix-only / removed (criterion-validity.csv): cognitive load -0.092 (0.6516) / +0.107
%     (0.5954); HR +0.088 (0.6774) / +0.036 (0.8647).
%   D-30: Round 1 participants were instructed to refrain from agitated movement because it
%     would corrupt the EEG; the instruction was NOT given in Round 2 -- the fifth asymmetry,
%     alongside capture software, axis convention, sampling rate and window.
% ============================================================================

The order in which the three films were shown was not the same for everyone, and the exception
is worth stating because it was deliberate. The first two participants watched the clinical
footage first, then the boring film, then the interesting one. After those two sessions the
order was fixed for everyone who followed---interesting, then clinical, then boring---on a
suspicion that the boring film was depressing whatever came after it, and the published paper
records that reasoning. The consequence is that the boring film ran last for ten of the twelve
participants and second for the remaining two.

Gaze deviation requires one further decision, because the measure is an angle and an angle needs
a center. Nothing in the data supplies one. Three candidates were computed for every
participant and every film. The first takes center to be the direction the headset itself was
pointing, so the measure asks how far the eye has turned from straight ahead.\footnote{This one
had to be reconstructed per round rather than fixed once, because the two capture programs used
different conventions for which axis points forward---Round 1 treating X as forward with Y
horizontal and Z vertical, Round 2 treating Z as forward with X horizontal and Y vertical.
A single hard-coded forward direction
would have given every Round 1 episode a meaningless offset near ninety degrees.} The second
takes center to be the middle direction of that particular recording, so the measure asks how
far the eye strayed from where it mostly sat during that film. The third takes center to be the
participant's own middle direction across all three of their films.

The third is the one scored against what participants reported, and the reason is the design
rather than the result. The whole analysis is within-subject: every value on both sides is
measured against that participant's own average across their three episodes, so each person
serves as their own control. Of the three centers, only the per-participant one is built that
same way. It removes what differs between people and has no bearing on the question---how the
headset happened to sit on a particular face---and it removes it in the domain where that
nuisance actually lives, by recentering the direction the eye is pointing before any angle is
taken. It preserves what differs within a person and does bear on the question: whether they
held their eye differently during one film than another. The per-recording center goes one step
too far, recentering inside each episode and erasing precisely the between-film difference the
analysis is looking for. The headset-forward center keeps that difference but carries the
fitting variation along with it. All three are reported, and no result appears without its
center named.\footnote{The ranking corroborates the reasoning without being its ground: against
reported boredom the per-participant center reaches +0.590, the per-recording center +0.524,
and the headset-forward center +0.359, which does not hold up. The two personal centers agree
in direction and in ordering, so nothing here turns on which of them is used.}

What each measurement can and cannot say follows from all of this, and is worth stating once
before presenting any of the results. Two measurements carry the analysis: gaze deviation, centered on the
participant, and eye closure. The self-report is the third, and it is the criterion the other
two are scored against rather than a competitor to them. Pupil diameter is reported and set
aside for reasons already given. Cognitive load was collected as a curiosity and was never a
candidate for analysis. The headset reports it as a single number derived by proprietary
calculation, and the calculation is not published: there is no way to establish which
variables enter it, how they are weighted, or how anything particular to this study would
affect what it returns. A measurement whose construction cannot be inspected cannot be
interpreted, whatever value it happens to report. Heart rate and its variability are removed for a different reason, and one that
belongs to the equipment rather than to the data. Both are taken by sensors that must rest
against the forehead, and the headset's housing is shaped to a narrow range of facial
structures; on a number of participants, most consistently those of Asian descent, it did not
sit flush against the skin. The contact those sensors require was therefore never established,
and no reading taken through a gap can be relied upon regardless of what it says. The EEG yields a direction and no result, as described. One asymmetry
between the rounds remains to be named alongside the capture software, the axis conventions,
the sampling rate and the analysis window: Round 1 participants were asked to refrain from
agitated movement, because such movement corrupts an EEG recording, and Round 2 participants,
who wore no EEG cap, were given no such instruction.
